\section{Classical ordinary control on the original function spaces}
\label{sec:ordinary-control-v42}

\status{This chapter proves the exact generic endpoint
\texttt{CtrlOrd.ctrlOrd}.  It uses the completed ordinary structure
chapter and writes out the classical interpolation, lattice, spectral,
and control arguments.  The source assumes a complete Noetherian local
\emph{domain} $O$ with finite residue field, an injective map into an
algebraically closed field $\Omega$, and a supplied weight representation
with all its original fields.  These hypotheses are retained.  In
particular this theorem is distinct from the ordinary base-change
theorem, which allows arbitrary CNL rings and nonflat local maps.  The
construction of an actual algebraic weight representation and the full
arithmetic specialization theorem remain separate producers.}

\subsection{The exact control problem}

Retain the original Frame $X$, groups $G=G_a\times G_p$, subgroups
$U_b,N,N_f,T,T_f,N^-(c),T_c(c)$, positive monoid $P$, and element $z$
from the ordinary structure chapter.  Let $(O,\mathfrak m)$ be a
complete Noetherian local domain with finite residue field, and let
$\iota:O\hookrightarrow\Omega$ be an injective ring map to an
algebraically closed field.  All occurrences of $O$ in $\Omega$ below
mean this specified map.  Fix $m\geq0$ and
\[
 \tau:U_b\longrightarrow\operatorname{GL}_m(O),\qquad
 \theta,\phi:T\longrightarrow O^\times,\qquad U\leq U_b.
\]
Assume the original $\mathrm{FinLevel}(U)$ and
$\mathrm{Free}(\tau,\theta,U)$ conditions, with their stabilizer
formula from OST1--OST3.  No module freeness is assumed here.

The supplied $\mathrm{WtRep}$ consists of an integer $d\geq0$,
\[
 \xi:G_p\longrightarrow\operatorname{GL}_{d+1}(\Omega),\qquad
 \chi:T_f\longrightarrow\Omega^\times,\qquad c_0\geq0,
\]
and exactly the following conditions.  Matrices described as integral
have entries in $\iota(O)$.
\begin{enumerate}
\item $\xi(n)$ for $n\in N$, $\xi(t)$ for $t\in T$, and
$\xi(y)$ for $y\in N^-(c_0)$ are integral.
\item $\chi(a)^{-1}\xi(a)$ is integral for every $a\in P$.
\item The zeroth row of $\xi(n)$ is $(1,0,\ldots,0)$ for
$n\in N_f$; the zeroth row of $\xi(t)$ is
$(\chi(t),0,\ldots,0)$ for $t\in T_f$.
\item Every entry outside column zero of
$\chi(z)^{-1}\xi(z)$ lies in $\iota(\mathfrak m)$.
\item For every $k\geq0$ there is $c_1$ such that
$\xi(y)-1$ has entries in $\iota(\mathfrak m^k)$ for
$y\in N^-(c_1)$.
\item For every $k\geq0$ there is $c_2$ such that
$\chi(t)-1\in\iota(\mathfrak m^k)$ for
$t\in T_f\cap T_c(c_2)$.
\item For $t\in T_f\cap Z(G_p)$, one has
$\xi(t)=\chi(t)\,1$.
\end{enumerate}
This list records all ten assertion fields of the original record;
the three integrality assertions in the first item are separate fields,
as are its two zeroth-row assertions.  Fix a classical depth $c$ with
\begin{equation}\label{eq:oct-depth}
 c\geq c_0,\qquad \phi|_{T_c(c)}=1,\qquad
 \iota(\theta(t))=\iota(\phi(t))\chi(t)\quad(t\in T).
\end{equation}

Use $\tau_\Omega$ and $\phi_\Omega$ for the original coefficientwise
images.  The actual classical space $C_c$ consists of families
$F=(F_i)_{0\leq i<m}$, $F_i:G\to\Omega^{d+1}$, satisfying
\begin{align}
 F_i(\operatorname{glob}(\gamma)g)
      &=\xi(\operatorname{glob}(\gamma)_p)F_i(g),\\
 \tau_\Omega(u)F(g(u,1))&=F(g)\quad(u\in U),\\
 F(g(1,n))&=F(g)\quad(n\in N),\\
 F(g(1,y))&=F(g)\quad(y\in N^-(c)),\\
 F(g(1,t))&=\iota(\phi(t))^{-1}F(g)\quad(t\in T).
 \label{eq:oct-classical}
\end{align}
In the away equation $\tau_\Omega$ acts on the $i$ index.  This is
literally $X.\mathrm{ClO}(W,\tau,U,c,\phi)$, with no replacement by
an abstract isomorphic space.

Define the untwisted family and control map by
\begin{equation}\label{eq:oct-control}
 \widetilde F_i(g)=\xi(g_p)^{-1}F_i(g),\qquad
 \operatorname{ctrl}(F)(g)_i=\widetilde F_i(g)_0.
\end{equation}
For $a\in P$ let $R_a$ be the same quotient representatives used by
the original $U_a$, and put
\begin{equation}\label{eq:oct-T}
 T_aF=\chi(a)^{-1}\sum_{n\in R_a}r_{na}F,
       \qquad T=T_z.
\end{equation}
On the full ambient space this denotes the source's fixed quotient
choices; on $N$-fixed forms its independence is OST5.  Write
\[
 \operatorname{Lat}(k)=
 \{F:\widetilde F_i(g)_j\in\iota(\mathfrak m^k)
                  \text{ for all }i,g,j\},\qquad
 \Lambda_c=C_c\cap\operatorname{Lat}(0).
\]
More generally $\operatorname{Lat}(D)$ has untwisted values in an
$O$-submodule $D\subset\Omega$.  The classical ordinary space is
the literal source sum
\begin{equation}\label{eq:oct-clord}
 C_c^{\mathrm{ord}}=
 \sum_{\substack{\mu\in\Omega^\times\\
                  \mu,\mu^{-1}\text{ integral over }O}}
           \left(C_c\cap\bigcup_{r\geq0}\ker(T-\mu)^r\right).
\end{equation}
Let $M=\mathrm{Sp}_O(\tau,\theta,U)$ and
$\mathcal O=\mathrm{Ord}_O(\tau,\theta,U)$ be the original spaces
of $O^m$-valued functions.  The coefficient map $f\mapsto\iota f$
is applied at every $g$ and every coordinate.

\subsection{Classical integral coordinates and the control formula}

% BEGIN OCT1
\paragraph{OCT1. Extending the finite character to the Iwahori subgroup.}
\label{claim:OCT1}
There is a character $\psi_c:I(c)\to O^\times$ that is trivial
on $N$ and $N^-(c)$ and equals $\phi$ on $T$.
\begin{proof}
Use the torus coordinate $t_0:I(0)\to T$ proved in OST2 and set
$\psi_c(b)=\phi(t_0(b))$.  If $b'=nty$ with
$n\in N,t\in T,y\in N^-(c)$, OST2 gives
\[
 \phi(t_0(bb'))=\phi(t_0(b))\phi(t),\qquad
 \phi(t_0(b'))=\phi(t).
\]
The possible extra torus factor lies in $T_c(c)$ and is killed by
$\phi$ by \eqref{eq:oct-depth}.  Thus $\psi_c(bb')=
\psi_c(b)\psi_c(b')$.  Taking only one of $n,t,y$ nonidentity
proves its three required restrictions.  This uses the actual
Iwahori torus coordinate, without postulating a torus quotient.
\end{proof}
% END OCT1

% BEGIN OCT2
\paragraph{OCT2. Interpolation of every untwisted classical value.}
\label{claim:OCT2}
Choose distinct finite double-coset representatives $x_1,\ldots,x_s$
at depth $c$ as in OST1.  For every list of matrices
$v_r\in\operatorname{Mat}_{m,d+1}(\Omega)$ there exists exactly
one $F\in C_c$ with $\widetilde F(x_r)=v_r$.
\begin{proof}
Let $K=U\times I(c)$ and define its action on such matrices by
\[
 \rho(u,b)v=\iota(\psi_c(b))\,
                   \tau_\Omega(u)v\xi(b)^{\mathsf T}.
\]
It is a representation: left multiplication, right multiplication
by the transposed matrix, and the scalar commute, and
$\xi(bb')^{\mathsf T}=\xi(b')^{\mathsf T}\xi(b)^{\mathsf T}$
gives the required composition order.

If $\operatorname{glob}(\gamma)x_r=x_r(u,b)$, original Free
says that $b\in Z(G_p)\cap T$ and
$\theta(b)\tau(u)=1$.  The central field of $W$ gives
$\xi(b)=\chi(b)1$, while OCT1 and \eqref{eq:oct-depth} give
\[
 \rho(u,b)v=\iota(\phi(b))\chi(b)\tau_\Omega(u)v
                       =\iota(\theta(b))\tau_\Omega(u)v=v.
\]
Thus every stabilizer acts trivially, on every prescribed matrix.
For $g=\operatorname{glob}(\gamma)x_r(u,b)$ define
\[
                 \Phi(g)=\rho((u,b)^{-1})v_r.
\]
The representative index is unique.  Comparing two factorizations
of the same $g$ yields a stabilizer relation and hence equality of
the two displayed values.  Therefore $\Phi$ is well defined, left
$\Gamma$-invariant, and satisfies
$\Phi(g(u,b))=\rho((u,b)^{-1})\Phi(g)$.

Now define $F_i(g)=\xi(g_p)\Phi_i(g)$, or in row-matrix notation
$F(g)=\Phi(g)\xi(g_p)^{\mathsf T}$.  The cancellation of
$\xi(b)^{-1}$ against $\xi(g_pb)$ shows
$F(g(1,b))=\iota(\psi_c(b))^{-1}F(g)$; the same formula for
$(u,1)$ gives the away equation.  Left invariance of $\Phi$ gives
the required twisted left action of $F$.  The three restrictions
of OCT1 give every local equation in \eqref{eq:oct-classical}.
Finally $\widetilde F=\Phi$, so its values are the prescribed
$v_r$.  Conversely any classical $F$ has precisely this
transformation law for $\widetilde F$ and is determined by those
values.  This proves uniqueness as well as existence.
\end{proof}
% END OCT2

% BEGIN OCT3
\paragraph{OCT3. The actual integral lattice is finite free.}
\label{claim:OCT3}
The map
\[
 E_c:C_c\longrightarrow\Omega^{s m(d+1)},\qquad
       E_c(F)_{r,i,j}=\widetilde F_i(x_r)_j
\]
is an $\Omega$-linear isomorphism.  For every $O$-submodule
$D\subset\Omega$,
\begin{equation}\label{eq:oct-lattice-test}
 F\in C_c\cap\operatorname{Lat}(D)\quad\Longleftrightarrow
          E_c(F)\in D^{s m(d+1)}.
\end{equation}
Consequently $\Lambda_c\simeq O^{s m(d+1)}$ and
$C_c\cap\operatorname{Lat}(k)=\mathfrak m^k\Lambda_c$.
\begin{proof}
The isomorphism is OCT2 and linearity of its defining formulas.
The forward implication of \eqref{eq:oct-lattice-test} is immediate.
For the converse write
$g=\operatorname{glob}(\gamma)x_r(u,b)$ and
$b=nty$ with $n\in N,t\in T,y\in N^-(c)$, by OST1.
Left multiplication by $\operatorname{glob}(\gamma)$ leaves
$\widetilde F$ unchanged.  Right multiplication by $(u,1)$ acts
on its $i$ index by $\tau_\Omega(u^{-1})$, an integral matrix.
For each local factor $b_0\in\{n,t,y\}$ one has
\[
 \widetilde F_i(h(1,b_0))=
          \iota(s_0)\xi(b_0^{-1})\widetilde F_i(h),
\]
where $s_0=1$ for $n,y$ and $s_0=\phi(t)^{-1}$ for $t$.
The three matrices are integral by the first fields of $W$;
for the last one use $N^-(c)\subset N^-(c_0)$.  Multiplication
by an integral matrix preserves $D$ in every coordinate, since
$D$ is an $O$-module.  Successively removing the factors proves
the reverse implication.  Taking $D=\iota(O)$ and then
$D=\iota(\mathfrak m^k)$ proves the two final assertions using
injectivity of $\iota$.  This is an assertion about the exact
value lattice of the original functions.
\end{proof}
% END OCT3

% BEGIN OCT4
\paragraph{OCT4. Exact equivariance of the original control map.}
\label{claim:OCT4}
The map in \eqref{eq:oct-control} is $\Omega$-linear and satisfies
\begin{equation}\label{eq:oct-equivariance}
 \operatorname{ctrl}(T_aF)=U_a\operatorname{ctrl}(F),\qquad
 \operatorname{ctrl}(\mathrm{raF}(u)F)
                         =\mathrm{ra}(u)\operatorname{ctrl}(F).
\end{equation}
Here the away actions use the same $\tau_\Omega$, and the first
identity holds for every ambient family and every $a\in P$.
\begin{proof}
Linearity follows immediately from multiplication by
$\xi(g_p)^{-1}$ and extraction of coordinate zero.  For
$n\in N_f,t\in T_f$ and any vector $v$, the zeroth-row fields
of $W$ give
\[
 \bigl(\xi(g_pnt)^{-1}v\bigr)_0
      =\chi(t)^{-1}\bigl(\xi(g_p)^{-1}v\bigr)_0.
\]
Indeed the inverse matrix factors as
$\xi(t^{-1})\xi(n^{-1})\xi(g_p)^{-1}$; the first two zeroth
rows successively multiply coordinate zero by $\chi(t)^{-1}$
and leave it unchanged.  Apply this with $t=a$ and the value
$v=F_i(g(1,na))$ for each original representative $n\in R_a$.
Summing is exactly the first identity.  The away action does
not change $g_p$, and its matrix acts on the separate $i$ index;
the two finite matrix sums commute, proving the second identity.
The same separation of factors proves that away actions commute
with every local right translation.  No invariance hypothesis
was needed for these identities.
\end{proof}
% END OCT4

% BEGIN OCT5
\paragraph{OCT5. Integral classical forms control to the actual ordinary ambient space.}
\label{claim:OCT5}
The character $\theta$ is continuous in the original sense.
Every $F\in\Lambda_c$ has a unique $f\in M$ with
$\iota f=\operatorname{ctrl}(F)$.  More precisely, if $c'\geq c$
and $\xi(y)-1$ has entries in $\iota(\mathfrak m^k)$ for
$y\in N^-(c')$, then
\begin{equation}\label{eq:oct-ctrl-continuity}
                r_yf-f\in\mathrm{valIn}_O(k)
                         \quad(y\in N^-(c')).
\end{equation}
\begin{proof}
Choose $c'\geq c$ with $\chi(t)-1\in\iota(\mathfrak m^k)$
on $T_c(c')$.  There $\phi(t)=1$, so
\eqref{eq:oct-depth} and injectivity of $\iota$ give
$\theta(t)-1\in\mathfrak m^k$.  This is ContChar.

Since $F$ is integral, all its untwisted values have unique
$O$-preimages.  Define $f(g)_i$ as the preimage of
$\widetilde F_i(g)_0$.  Twisted left equivariance cancels in
$\widetilde F$, giving the left invariance of $f$.  The zeroth
row of $\xi(n^{-1})$ proves its $N$-invariance.  The zeroth
row of $\xi(t^{-1})$, together with
$\iota(\theta(t))=\iota(\phi(t))\chi(t)$, gives its
$\theta^{-1}$ torus law.  The away equation follows from OCT4
and injectivity of $\iota$.

For continuity take $y\in N^-(c')$.  Since $c'\geq c$,
$F(g(1,y))=F(g)$, and therefore
\[
 \widetilde F_i(g(1,y))=\xi(y^{-1})\widetilde F_i(g).
\]
The inverse $y^{-1}$ lies in the same subgroup.  If
$\widetilde F_i(g)_j=\iota(v_j)$ and
$\xi(y^{-1})_{0j}=\delta_{0j}+\iota(a_j)$ with
$a_j\in\mathfrak m^k$, then
\[
                  f(g(1,y))_i-f(g)_i=\sum_j a_jv_j
                                                  \in\mathfrak m^k.
\]
This proves \eqref{eq:oct-ctrl-continuity}, uniformly in $g,i$.
The continuity field of $W$ supplies such a depth for every $k$,
so $f$ satisfies every condition of the original $M$.  Its
uniqueness and the $O$-linearity of $F\mapsto f$ follow from
injectivity of $\iota$.
\end{proof}
% END OCT5

% BEGIN OCT6
\paragraph{OCT6. Integral stability and strict contraction in the control kernel.}
\label{claim:OCT6}
Every $T_a$, $a\in P$, preserves $\operatorname{Lat}(D)$ for
every $O$-submodule $D\subset\Omega$.  More specifically, if
$\mathfrak m D\subset D'$, $F\in\operatorname{Lat}(D)$, and
all control values of $F$ lie in $D'$, then
\begin{equation}\label{eq:oct-contraction}
                         TF\in\operatorname{Lat}(D').
\end{equation}
\begin{proof}
The untwisted value of one normalized summand in $T_aF$ is
\begin{equation}\label{eq:oct-summand}
 \xi(n)\bigl(\chi(a)^{-1}\xi(a)\bigr)
                         \widetilde F_i(g(1,na)).
\end{equation}
This follows by substituting
$F_i(h)=\xi(h_p)\widetilde F_i(h)$ in the defining sum and
cancelling $\xi(g_p)^{-1}\xi(g_p)$.  Both matrices in
\eqref{eq:oct-summand} are integral by the fields of $W$, so
each summand preserves $D$.

For $a=z$, every column of the second matrix except column
zero has entries in $\iota(\mathfrak m)$.  Their contributions
therefore lie in $D'$.  The contribution of column zero lies
in $D'$ because its input is the control value and its
coefficient is integral.  Multiplication by $\xi(n)$ preserves
$D'$.  The finite sum proves \eqref{eq:oct-contraction}.
In particular $T$ preserves every $\operatorname{Lat}(k)$, and
if $\operatorname{ctrl}(F)=0$ and
$F\in\operatorname{Lat}(k)$, then
$TF\in\operatorname{Lat}(k+1)$.
\end{proof}
% END OCT6

% BEGIN OCT7
\paragraph{OCT7. Classical level lowering on the original functions.}
\label{claim:OCT7}
For $c'\geq c$ one has $C_c\subset C_{c'}$.  Every $T_a$
preserves $C_c$ and $\Lambda_c$.  For integers $e,n\geq0$
with $n\geq e$,
\begin{equation}\label{eq:oct-lowering}
              T^n C_{c+e}\subset C_c,\qquad
              T^n\Lambda_{c+e}\subset\Lambda_c.
\end{equation}
\begin{proof}
The inclusion of levels follows from
$N^-(c')\subset N^-(c)$.  Twisted left equivariance is
preserved term by term under right translation, and away
equivariance commutes with it by OCT4.  For local invariance,
the permutation-of-cosets proof OST6 applies to the
representation on the classical families: $\phi(T_c(c))=1$
gives its required torus-congruence invariance.  The same
proof gives preservation of the $N$ and torus character laws
by OST5.  Multiplication by the scalar $\chi(a)^{-1}$ does
not alter any of these laws.  Thus $T_aC_c\subset C_c$;
OCT6 adds lattice preservation.

At each depth $c+r$, $\phi$ remains trivial on $T_c(c+r)$.
The strict part of OST6 gives
$T C_{c+r+1}\subset C_{c+r}$.  Apply it $e$ times and then
use stability at depth $c$ for the remaining $n-e$ iterates.
OCT6 proves the integral assertion at every step.  The
normalizing factor remains the original $\chi(z)^{-1}$.
\end{proof}
% END OCT7

\subsection{A finite lattice and its actual spectral ordinary part}

% BEGIN OCT8
\paragraph{OCT8. Factorial decomposition of a finite integral matrix.}
\label{claim:OCT8}
Let $L\simeq O^r$ be finite free and $T_0\in\operatorname{End}_O(L)$.
There is one idempotent $e_0$ commuting with $T_0$ such that
\[
 e_0=\lim_{j\to\infty}T_0^{j!},\qquad
 L=S_0\oplus Q_0,\quad S_0=e_0L,
                         \quad Q_0=(1-e_0)L.
\]
Both summands are finite free.  The restriction of $T_0$ to
$S_0$ is invertible; for every $k$ some $n$ satisfies
$T_0^nQ_0\subset\mathfrak m^k Q_0$.  Moreover
\begin{equation}\label{eq:oct-periodic}
 v\in S_0\quad\Longleftrightarrow\quad
 \forall k\ \exists q\geq1\quad T_0^qv-v\in\mathfrak m^kL.
\end{equation}
At every $k$ the factorial convergence is uniform in $v\in L$.
\begin{proof}
Each $O/\mathfrak m^k$ is finite.  Indeed
$\mathfrak m$ is finitely generated, so each
$\mathfrak m^j/\mathfrak m^{j+1}$ is finite-dimensional over
the finite residue field; induction along these quotients
gives finiteness.  Thus the endomorphism monoid of
$L/\mathfrak m^kL$ is finite.  Some $b\geq0,q\geq1$ satisfy
$T_0^{n+q}=T_0^n$ on this quotient for every $n\geq b$.
For all sufficiently large $j$, $j!\geq b$ and $q\mid j!$;
the endomorphism $T_0^{j!}$ is then one fixed idempotent $e_k$.
Idempotence follows by comparing the exponents $2j!$ and $j!$.
On its image $T_0^{j!}$ is the identity, so $T_0$ is invertible
there; on its kernel that power is zero.

These $e_k$ are compatible, since every one is given by all
sufficiently late members of the same factorial sequence.
Completeness of $O$ constructs their matrix-entry limit $e_0$;
separatedness makes it unique and gives $e_0^2=e_0$ and
$e_0T_0=T_0e_0$.  The decomposition is now literal.  Its
summands are finite projective and hence finite free over the
local ring; the elementary projective-to-free argument is
written in OST17.  Their reductions identify with the image
and kernel of $e_k$, because the decomposition is split.
On a basis of $S_0$, the determinant of $T_0$ is a unit modulo
$\mathfrak m$, hence a unit of $O$; the adjugate supplies the
inverse over $O$.  The kernel statement modulo each
$\mathfrak m^k$ gives the asserted topological nilpotence.

If the right side of \eqref{eq:oct-periodic} holds, take a
factorial divisible by its period $q$ at precision $k$.
Then $e_kv=v$ modulo $\mathfrak m^k$, for every $k$, whence
$e_0v=v$.  Conversely $T_0$ acts as an automorphism of the
finite module $S_0/\mathfrak m^kS_0$, so some positive power
fixes the given vector.  The construction of $e_k$ as an
endomorphism, rather than as a limit for one vector, proves
the asserted uniform convergence.  The case $r=0$ satisfies
the same statements with zero modules.
\end{proof}
% END OCT8

% BEGIN OCT9
\paragraph{OCT9. The finite lattice ordinary part equals the integral-unit spectrum.}
\label{claim:OCT9}
Extend the data of OCT8 to $L_\Omega=\Omega\otimes_O L$.
The $\Omega$-span of $S_0$ is exactly the sum of generalized
eigenspaces of $T_0$ whose eigenvalues $\mu$ satisfy
$\mu\ne0$ and both $\mu,\mu^{-1}$ integral over $O$.
\begin{proof}
On $S_0$, both $T_0$ and $T_0^{-1}$ are integral matrices.
Their monic characteristic polynomials show that every
eigenvalue $\mu$ and its inverse are integral over $O$;
invertibility shows $\mu\ne0$.

On $Q_0$, choose $n\geq1$ with
$T_0^nQ_0\subset\mathfrak m Q_0$.  If $q$ is its rank,
the characteristic polynomial of $T_0^n|_{Q_0}$ is
\[
             X^q+a_1X^{q-1}+\cdots+a_q,
                              \qquad a_j\in\mathfrak m^j.
\]
The coefficient condition follows by expanding principal
minors: each degree-$j$ minor is a sum of products of $j$
entries in $\mathfrak m$.  Suppose an eigenvalue $\mu$ of
$T_0|_{Q_0\otimes\Omega}$ were nonzero with $\mu^{-1}$
integral over $O$.  Substitution of $\mu^n$ and division by
$\mu^{nq}$ would give
\[
              1+a_1\mu^{-n}+\cdots+a_q\mu^{-nq}=0.
\]
The algebra $B=O[\mu^{-1}]\subset\Omega$ is a finite
$O$-module, since $\mu^{-1}$ satisfies a monic equation.
The displayed identity gives $1\in\mathfrak m B$, so
$B=\mathfrak m B$.  Nakayama gives $B=0$, contradicting
$1\ne0$ in $\Omega$.  Thus the complementary block has no
such eigenvalue.  If $q=0$ there is no complementary block.

For completeness, over an algebraically closed field the
characteristic polynomial factors as
$\prod_\mu(X-\mu)^{r_\mu}$.  The factors are pairwise
coprime; Bezout identities in $\Omega[X]$, evaluated at the
operator and combined with its characteristic equation,
decompose the vector space into the kernels
$\ker(T_0-\mu)^{r_\mu}$.  These are precisely the maximal
generalized eigenspaces.  Apply this decomposition separately
to the two $T_0$-stable summands of OCT8.  All eigenvalues on
the first block, and none on the second, meet the required
integral-unit condition.  This proves the equality.  No
valuation on $\Omega$ or normality of $O$ was used.
\end{proof}
% END OCT9

% BEGIN OCT10
\paragraph{OCT10. A finite free lattice spanning the exact classical ordinary space.}
\label{claim:OCT10}
For the original $C_c,\Lambda_c,T$, let $e_c$ be the factorial
projector on $\Lambda_c$ supplied by OCT8 and set
$S_c=e_c\Lambda_c$.  Then $S_c$ is finite free,
\begin{equation}\label{eq:oct-spectral}
                      C_c^{\mathrm{ord}}=\Omega S_c,
\end{equation}
and every $F\in S_c$ is periodic modulo every
$\operatorname{Lat}(k)$.  For every $G\in\Lambda_c$ and
every $k$,
\begin{equation}\label{eq:oct-classical-factorial}
 \exists J\ \forall j\geq J\quad
              T^{j!}G-e_cG\in\operatorname{Lat}(k).
\end{equation}
The choice of $e_c$ precedes $G,k$, and its threshold can be
chosen uniformly in $G\in\Lambda_c$ for each $k$.
\begin{proof}
By OCT3 the lattice is finite free and its extension to
$\Omega$ is exactly $C_c$ through the evaluation isomorphism.
OCT7 says that $T$ is an $O$-linear endomorphism of this
lattice.  Apply OCT8, using OCT3 to identify its intrinsic
filtration $\mathfrak m^k\Lambda_c$ with its global
untwisted-value filtration.  OCT9 gives
\eqref{eq:oct-spectral}.  Because $T$ preserves $C_c$,
the generalized eigenspace of its restriction at $\mu$ is
exactly $C_c\cap\bigcup_r\ker(T-\mu)^r$ in the original
ambient function space; the defining iterates are the same
on $C_c$.  Thus the spectral sum is exactly
\eqref{eq:oct-clord}, not a substituted definition.
\end{proof}
% END OCT10

\subsection{Integral control and its inverse}

% BEGIN OCT11
\paragraph{OCT11. Periodic integral elements in the control kernel vanish.}
\label{claim:OCT11}
Suppose an ambient family $F$ lies in $\operatorname{Lat}(0)$,
has $\operatorname{ctrl}(F)=0$, and for every $k$ has a
positive period $q$ modulo $\operatorname{Lat}(k)$:
$T^qF-F\in\operatorname{Lat}(k)$.  Then $F=0$.
\begin{proof}
Induct on $k$ to show $F\in\operatorname{Lat}(k)$.  The
initial assertion is assumed.  If it holds at $k$, OCT6 gives
$TF\in\operatorname{Lat}(k+1)$, and stability in OCT6 gives
$T^qF\in\operatorname{Lat}(k+1)$ for every $q\geq1$.
Choose a positive period at precision $k+1$ and subtract its
congruence to obtain $F\in\operatorname{Lat}(k+1)$.

For each untwisted coordinate, choose its $O$-preimage using
membership in $\operatorname{Lat}(0)$.  Injectivity of
$\iota$ and membership at every higher precision place that
same preimage in $\bigcap_k\mathfrak m^k=0$.  Therefore all
untwisted coordinates vanish.  Multiplying by the invertible
matrix $\xi(g_p)$ gives $F_i(g)=0$ for all $i,g$.
This argument does not presume that $F$ belongs to a finite
dimensional space, and uses positive, rather than zero,
periods.
\end{proof}
% END OCT11

% BEGIN OCT12
\paragraph{OCT12. The integral ordinary lattice maps injectively to the original ordinary module.}
\label{claim:OCT12}
There is a unique $O$-linear map
\[
             q_c:S_c\longrightarrow\mathcal O,
                  \qquad \iota(q_cF)=\operatorname{ctrl}(F),
\]
and it is injective.
\begin{proof}
OCT5 first supplies the unique $f\in M$.  At any precision
$k$, OCT10 gives $T^qF-F\in\operatorname{Lat}(k)$ for some
$q\geq1$.  Its zeroth untwisted coordinates, OCT4, and
coefficientwise naturality of the original finite sums give
\[
                     U_z^q f-f\in\mathrm{valIn}_O(k).
\]
Here membership follows from injectivity of $\iota$:
an $O$-value whose image lies in $\iota(\mathfrak m^k)$
already lies in $\mathfrak m^k$.  The literal periodicity
characterization OST10 therefore gives $f\in\mathcal O$.
Linearity follows by uniqueness.  If $q_cF=0$, OCT10 supplies
all the periods needed in OCT11, which then gives $F=0$.
\end{proof}
% END OCT12

% BEGIN OCT13
\paragraph{OCT13. An integral classical approximation at the required fixed depth.}
\label{claim:OCT13}
For every $f\in\mathcal O$ there are
$G\in\Lambda_c$ and $h\in M$ such that
\begin{equation}\label{eq:oct-approx}
 \iota h=\operatorname{ctrl}(G),\qquad
                         h-f\in\mathrm{valIn}_O(1).
\end{equation}
\begin{proof}
Choose a depth at which $f$ is $N^-$-invariant modulo
$\mathfrak m$, and a depth at which $\xi(y)-1$ has all
entries in $\iota(\mathfrak m)$.  Take $c'$ at least both
these depths and $c$, and write $e=c'-c$.  Triviality of
$\phi$ on $T_c(c')$ still holds.  Choose the finite distinct
representatives $x_r$ at this deeper depth.  OCT2 constructs
$F'\in C_{c'}$ with
\[
      \widetilde F'_i(x_r)_j=
          \begin{cases}\iota(f(x_r)_i),&j=0,\\0,&j\ne0.
          \end{cases}
\]
These are integral values, so OCT3 gives $F'\in\Lambda_{c'}$.
OCT5 gives its integral control $h'\in M$, with
$h'(x_r)=f(x_r)$.  Its refined continuity assertion says
that $h'$ is $N^-(c')$-invariant modulo $\mathfrak m$.
The same is true for $f$ by the choice of $c'$.  Apply
OST4 to $h'-f$, whose sampled values are zero, to obtain
$h'-f\in\mathrm{valIn}_O(1)$.

Choose a period $q\geq1$ for $f$ at this precision, and a
multiple $M_0$ of $q$ with $M_0\geq e$, for example
$M_0=q(e+1)$.  Set
\[
                         G=T^{M_0}F',\qquad h=U_z^{M_0}h'.
\]
OCT7 lowers $G$ from depth $c'$ to depth $c$, and OCT6
preserves its integrality.  OST7 gives $h\in M$.
OCT4 and coefficientwise naturality give
$\iota h=\operatorname{ctrl}(G)$.  Finally
\[
 h-f=U_z^{M_0}(h'-f)+(U_z^{M_0}f-f)
                                             \in\mathrm{valIn}_O(1).
\]
The first term belongs to the value filtration by OST5.
For the second, a positive multiple of a period is a
period: expand
$U_z^{qt}f-f=\sum_{a=0}^{t-1}U_z^{qa}(U_z^qf-f)$.
This proves \eqref{eq:oct-approx}.  The deeper depth is an
intermediate approximation choice; the final classical form
is at the original prescribed depth $c$.
\end{proof}
% END OCT13

% BEGIN OCT14
\paragraph{OCT14. Factorial projection gives surjectivity modulo the maximal ideal.}
\label{claim:OCT14}
The image of $q_c$ satisfies
\[
                         q_c(S_c)+\mathfrak m\mathcal O
                                                   =\mathcal O.
\]
\begin{proof}
For $f\in\mathcal O$, choose $G,h$ as in OCT13.  Put
$F_0=e_cG\in S_c$ and $f_0=q_cF_0\in\mathcal O$.
By OCT10, at precision one there is $J$ such that
$T^{j!}G-F_0\in\operatorname{Lat}(1)$ for $j\geq J$.
Choose a positive period $q$ for $f$ modulo
$\mathrm{valIn}_O(1)$ and then choose $j\geq\max(J,q)$.
The factorial $j!$ is divisible by $q$.  Applying control
and using OCT4 gives
\[
 U_z^{j!}h-f_0\in\mathrm{valIn}_O(1),\qquad
 U_z^{j!}f-f\in\mathrm{valIn}_O(1).
\]
Also $U_z^{j!}(h-f)$ lies in that filtration by OST5.
The identity
\[
 f-f_0=(U_z^{j!}h-f_0)-U_z^{j!}(h-f)-(U_z^{j!}f-f)
\]
therefore places $f-f_0$ in $\mathrm{valIn}_O(1)$.
Both $f$ and $f_0$ are ordinary, so the exact reverse
filtration identity OST16 gives
$f-f_0\in\mathfrak m\mathcal O$.
This step uses the established intrinsic filtration on
$\mathcal O$; a value congruence by itself would not suffice.
\end{proof}
% END OCT14

% BEGIN OCT15
\paragraph{OCT15. Integral control is an isomorphism.}
\label{claim:OCT15}
The map $q_c:S_c\to\mathcal O$ is an $O$-linear isomorphism.
\begin{proof}
OCT5 gives ContChar, so OST19 applies with the original
FinLevel and Free hypotheses and makes $\mathcal O$ finite
free.  The quotient $Q=\mathcal O/q_c(S_c)$ is therefore a
finite $O$-module.  OCT14 gives $Q=\mathfrak m Q$.
Nakayama gives $Q=0$, so $q_c$ is surjective.  Its injectivity
was proved in OCT12.  One may see the used form of Nakayama
directly: if generators $v_1,\ldots,v_r$ of $Q$ express each
$v_i$ as $\sum_j a_{ij}v_j$ with $a_{ij}\in\mathfrak m$,
then the adjugate of $1-(a_{ij})$ shows that its determinant
annihilates every generator.  This determinant lies in
$1+\mathfrak m$ and is a unit, forcing $Q=0$.
Thus the inverse lifts ordinary functions to actual classical
families and uses exactly the control formula
\eqref{eq:oct-control}.
\end{proof}
% END OCT15

% BEGIN OCT16
\paragraph{OCT16. Independent integral functions remain independent over the ambient field.}
\label{claim:OCT16}
Let $f_1,\ldots,f_r$ be $O$-linearly independent functions
from any set $D$ to $O$.  Then their coefficientwise images
are $\Omega$-linearly independent in $\Omega^D$.
\begin{proof}
The embedding of $O$ extends to an embedding of its fraction
field $K$ into $\Omega$.  Clearing finitely many denominators
in a $K$-linear relation shows that the $f_i$, now viewed as
$K$-valued functions, remain independent.  Let $V$ be their
$r$-dimensional $K$-span.  Evaluation at the points of $D$
spans $V^*$: otherwise the common kernel of these linear
functionals would contain a nonzero member of $V$, a
function zero at every point, which is impossible.  Choose
$r$ evaluations forming a basis of $V^*$.  The resulting
$r$-by-$r$ matrix $(f_i(x_j))$ has nonzero determinant in
$K$, and its entries and determinant lie in $O$.
Injectivity of $\iota$ makes that determinant nonzero in
$\Omega$.  Evaluating any $\Omega$-linear relation at those
$r$ points now forces every coefficient to vanish.  For
$r=0$ the statement is vacuous.  This is the precise point
where the source's domain hypothesis is used; it is not
silently removed from the control theorem.
\end{proof}
% END OCT16

% BEGIN OCT17
\paragraph{OCT17. The exact original classical ordinary control endpoint.}
\label{claim:OCT17}
Under exactly the standing hypotheses and
\eqref{eq:oct-depth}, the source's original map satisfies
\begin{align}
 F\in C_c^{\mathrm{ord}},\quad\operatorname{ctrl}(F)=0
                                     &\Longrightarrow F=0,\\
 \operatorname{ctrl}(C_c^{\mathrm{ord}})
                   &=\operatorname{span}_{\Omega}
                                \{\iota f:f\in\mathcal O\}.
 \label{eq:oct-endpoint}
\end{align}
These are precisely the two conclusions of
\texttt{UpperTheory.V.Hi.ctrlOrd}.
\begin{proof}
By OCT10, $C_c^{\mathrm{ord}}=\Omega S_c$.
Take an $O$-basis $F_1,\ldots,F_r$ of $S_c$ and put
$f_i=q_cF_i$.  By OCT15 these are an $O$-basis of the
original $\mathcal O$.  Apply OCT16 to their coordinates
as functions on $D=G\times\{0,\ldots,m-1\}$; the resulting
$\Omega$-valued functions $\iota f_i=\operatorname{ctrl}(F_i)$
are independent.  Every element of $\Omega S_c$ is a finite
$\Omega$-linear combination of the $F_i$, so zero control
forces all its coefficients to vanish.  Linearity and the
same basis show that its image is exactly the displayed
span.  This also shows that the coefficientwise tensor map
$\Omega\otimes_O\mathcal O\to\Omega^{G\times\{0,\ldots,m-1\}}$
is injective, and control identifies $C_c^{\mathrm{ord}}$
with its image.  All spaces in this argument are the
original spaces of functions from the source definitions.
\end{proof}
% END OCT17

\subsection{Compatibility and the remaining installation boundary}

% BEGIN OCT18
\paragraph{OCT18. Compatibility with original positive operators and classical depth.}
\label{claim:OCT18}
Every $T_a$, $a\in P$, preserves $S_c$, and
\[
                  q_cT_a=U_aq_c.
\]
It is bijective on $S_c$.  If $c'\geq c$ satisfies the same
standing depth requirements, then $S_c=S_{c'}$ as submodules
of the original ambient families, and their integral control
inverses are identical.  Thus the ordinary classical subspace
and the inverse control map are independent of sufficiently
deep classical level, with no new choices at each precision.
\begin{proof}
OCT7 and OCT6 show that $T_a$ preserves $\Lambda_c$.
On $N$-fixed families, the finite double-coset product rule
OST5 and multiplicativity of $\chi$ give
$T_aT=T T_a$.  Consequently $T_a$ preserves the periodic
submodule $S_c$ characterized by OCT8, or equivalently
commutes with its factorial projector.  OCT4 gives
equivariance under $q_c$.  By OST20 the original $U_a$ is
bijective on $\mathcal O$, and OCT15 therefore makes $T_a$
bijective on $S_c$.

For depths, OCT7 gives $\Lambda_c\subset\Lambda_{c'}$,
and periodicity in the global value lattices gives
$S_c\subset S_{c'}$.  The two maps $q_c,q_{c'}$ are
restrictions of the same pointwise formula and are both
isomorphisms onto $\mathcal O$.  For $F'\in S_{c'}$ choose
$F\in S_c$ with $q_cF=q_{c'}F'$.  Their difference belongs
to $S_{c'}$ and has zero control, so OCT12 gives $F'=F$.
This proves equality and identity of the inverse maps.
OCT10 then also gives equality of the classical ordinary
$\Omega$-spaces at these depths.
\end{proof}
% END OCT18

% BEGIN OCT19
\paragraph{OCT19. Compatibility with finite sums of actual away translations.}
\label{claim:OCT19}
Suppose two level/depth pairs $(U,c)$ and $(U',c')$ satisfy
the standing hypotheses for the same $O,\Omega,W,\tau,
\theta,\phi$.  Let
\[
       A_F=\sum_{u\in E}\mathrm{raF}(u),\qquad
       A_f=\sum_{u\in E}\mathrm{ra}(u),
\]
where $E$ is any finite indexed list of elements of $U_b$,
with repetitions retained.  If $A_F$ maps $C_c(U)$ into
$C_{c'}(U')$, then it maps $S_c(U)$ into $S_{c'}(U')$ and
\[
                         q_{c',U'}A_F=A_f q_{c,U}.
\]
In particular the same assertion holds for the natural
inclusion whenever $U'\subset U$ and $c'\geq c$.
\begin{proof}
An away translation acts on untwisted coordinates by the
integral matrix $\tau_\Omega(u)$ and an away translation
of $g$.  It preserves every global lattice
$\operatorname{Lat}(k)$, as does a finite sum.  It also
commutes with $T$, since the away and local actions commute
by OCT4.  Under the assumed classical-space inclusion, a
periodic element of $\Lambda_c(U)$ therefore maps to an
integral periodic element of $C_{c'}(U')$.  The criterion
OCT8 and OCT3 put its image in $S_{c'}(U')$.
OCT4 gives the displayed equality after applying $\iota$;
injectivity of $\iota$ gives the equality over $O$.
For natural inclusion use the singleton list $u=1$ and
the defining equations for the two levels.  Thus these
compatibilities concern the original translations and
actual inclusions, rather than an abstract identification
of modules of the same rank.
\end{proof}
% END OCT19

\paragraph{Exact source and proof boundary.}
The source pin is
\begin{center}
\texttt{7cbf7a0c32fca1badf05931879a37c3f8a2b03b4}.
\end{center}
The source definitions are
\texttt{ContinuousTorusCharacter.lean} (the classical section)
and \texttt{CtrlDefs.lean}; the endpoint is
\texttt{OrdinaryHeckeCharacters/CtrlOrd.lean}, with the control
identities in \texttt{CtrlEq.lean}.  The interpolation,
integral-lattice, spectral, and continuity helper files and
their exact byte hashes are recorded in the accompanying
source manifest.  OCT1--OCT19 are written proofs; the source
locations identify what they reconstruct and are not used
as black-box justifications of those proofs.

The endpoint quantifies over every supplied original Frame
and every supplied $W$ satisfying all ten assertion
fields, every $m$, every $\tau,\theta,\phi$, and every
admissible $U,c$.  It proves the exact domain-valued control
theorem; it neither changes the source to arbitrary
non-domain coefficients nor discards the specified embedding
$O\hookrightarrow\Omega$.  The arithmetic use in
\texttt{SpecIntIff.spec} first supplies such a $W$ over a
finite coefficient extension, a character $\phi$, and the
compatibility \eqref{eq:oct-depth}.  That construction,
the arithmetic-fiber norm dictionary, the full single-map
specialization theorem, and the same-ideal Hida/automorphy
endpoints are not proved merely by invoking OCT17.  They
remain separate obligations unless independently supplied
elsewhere in the assembled reconstruction.  No whole
\texttt{COUNT\_II\_SUPPORT} leaf is retired by this chapter.
